\SixthArticle{45}

Sei $k$ eine positive ganze Zahl (oder auch $=0$), und bezeichnen wir die ersten $k+1$ Glieder der Reihe $F(\alpha,\beta,\gamma,x)$ mit $X$. Das nächste Glied ist dann
\[
=\frac{\alpha.\alpha+1.\alpha+2\ .\ .\ .\ .\ \alpha+k.\beta.\beta+1.\beta+2\ .\ .\ .\ .\ \beta+k}
{1\ .\ 2\ .\ 3\ .\ .\ .\ .\ k+1.\gamma.\gamma+1.\gamma+2\ .\ .\ .\ .\ \gamma+k}x^{k+1}
\]
das sich auch so darstellen lässt:
\[
\frac{\Pi(\gamma-1)}{\Pi(\alpha-1)\Pi(\beta-1)}\ .\
\frac{\Pi(\alpha+k)\Pi(\beta+k)}{\Pi(k+1)\Pi(\gamma+k)}x^{k+1}
\]
und ebenso die folgenden Glieder. Daraus folgt:
\[
\mathrm{I.}\qquad
\frac{\Pi(\alpha+\beta-\gamma)(\Pi(-\gamma)}{\Pi(\alpha-\gamma)\Pi(\beta-\gamma)}\ .\ F(\alpha,\beta,\gamma,x)
\quad\text{lässt sich ausdrücken durch}
\]
\[
\begin{aligned}&\frac{\Pi(\alpha+\beta-\gamma)\Pi(-\gamma)}{\Pi(\alpha-\gamma)\Pi(\beta-\gamma)}\ .\ X\\&+\frac{\Pi(\alpha+\beta-\gamma)\Pi(-\gamma)\Pi(\gamma-1)}{\Pi(\alpha-1)\Pi(\beta-1)\Pi(\alpha-\gamma)\Pi(\beta-\gamma)}\\&\qquad\Sigma\left\{\frac{\Pi(\alpha+k+t)\Pi(\beta+k+t)}{\Pi(k+t+1)\Pi(\gamma+k+t)}x^{k+1+t}\right\}\end{aligned}
\]
wenn man sich für $t$ alle Werte $0,1,2,3$ usw. bis ins Unendliche eingesetzt denkt.

Ebenso lässt sich $F(\alpha+1-\gamma,\beta+1-\gamma,2-\gamma,x)$ ausdrücken durch
\[
\frac{\Pi(1-\gamma)}{\Pi(\alpha-\gamma)\Pi(\beta-\gamma)}
\Sigma\left(\frac{\Pi(\alpha-\gamma+t)\Pi(\beta-\gamma+t)}{\Pi t\,\Pi(1-\gamma+t)}x^t\right)
\]
wobei $t$ wie zuvor bestimmt wird. Da nun die Beziehungen $\Pi(1-\gamma)=(1-\gamma)\Pi(-\gamma)$ und $\Pi(\gamma-1)=-(1-\gamma)\Pi(\gamma-2)$ gelten, gilt offenbar
\SourceLeaf{224}
\[
\begin{aligned}[c]
\mathrm{II.}\qquad&\frac{\Pi(\alpha+\beta-\gamma)\Pi(\gamma-2)}{\Pi(\alpha-1)\Pi(\beta-1)}x^{1-\gamma}F(\alpha+1-\gamma,\beta+1-\gamma,2-\gamma,x)\\
&=-\frac{\Pi(\alpha+\beta-\gamma)\Pi(-\gamma)\Pi(\gamma-1)}{\Pi(\alpha-1)\Pi(\beta-1)\Pi(\alpha-\gamma)\Pi(\beta-\gamma)}
\\&\qquad\Sigma\left\{\frac{\Pi(\alpha-\gamma+t)\Pi(\beta-\gamma+t)}{\Pi t\,\Pi(1-\gamma+t)}x^{1+t-\gamma}\right\}
\end{aligned}
\]
Daher lässt sich Formel 86 auch folgendermaßen darstellen:
\[
\begin{aligned}[c]
&F(\alpha,\beta,\alpha+\beta+1-\gamma,1-x)\\
&=\frac{\Pi(\alpha+\beta-\gamma)\Pi(-\gamma)}{\Pi(\alpha-\gamma)\Pi(\beta-\gamma)}X\\
&\quad+\frac{\Pi(\alpha+\beta-\gamma)\Pi(-\gamma)\Pi(\gamma-1)}{\Pi(\alpha-1)\Pi(\beta-1)\Pi(\alpha-\gamma)\Pi(\beta-\gamma)}
\\&\qquad\Sigma\Bigl\{\frac{\Pi(\alpha+k+t)\Pi(\beta+k+t)}{\Pi(k+t+1)\Pi(\gamma+k+t)}x^{k+1+t}
\\&\qquad-\frac{\Pi(\alpha-\gamma+t)\Pi(\beta-\gamma+t)}{\Pi(t-\gamma+1)\Pi t}x^{1+t-\gamma}\Bigr\}
\end{aligned}
\]
Dieser Ausdruck zeigt unmittelbar, dass jede der Differenzen unter dem Zeichen $\Sigma$ $=0$ wird, wenn man $\gamma=-k$ setzt. Da aber zugleich $\Pi(\gamma-1)$ hier unendlich groß wird, kann das Produkt offenbar endlich sein. Um seinen Wert durch endliche Größen auszudrücken, setzen wir zunächst $\gamma+k=\omega$; daraus folgt
\[
\Pi(\gamma-1).\gamma.(\gamma+1)(\gamma+2)\ .\ .\ .\ .\ (\gamma+k-1)\omega=\Pi\omega
\]
oder
\[
\Pi(\gamma-1)=\frac{\Pi\omega}{\omega(\omega-1)(\omega-2)\ .\ .\ .\ .\ (\omega-k)}
\]
Es kommt also im Wesentlichen darauf an, dass wir sehen, was aus
\[
\begin{aligned}&\frac{1}{\omega}\Bigl\{\frac{\Pi(\alpha-\gamma+t+\omega)\Pi(\beta-\gamma+t+\omega)}{\Pi(t-\gamma+1+\omega)\Pi(t+\omega)}x^{1+t-\gamma+\omega}\\&\qquad-\frac{\Pi(\alpha-\gamma+t)\Pi(\beta-\gamma+t)}{\Pi(t-\gamma+1)\Pi t}x^{1+t-\gamma}\Bigr\}\end{aligned}
\]
wird, wenn $\omega$ gegen null abnimmt. Nach bekannten Grundsätzen ergibt sich daraus jedoch
\[
-\frac{\mathrm{d}U}{\mathrm{d}\gamma}
\]
wenn wir der Kürze halber
\[
\frac{\Pi(\alpha-\gamma+t)\Pi(\beta-\gamma+t)}{\Pi(t-\gamma+1)\Pi(t-k-\gamma)}x^{1+t-\gamma}=U
\]
setzen und allein $\gamma$ als veränderlich betrachten. Daraus folgt aber
\[
\begin{aligned}&\frac{\mathrm{d}U}{U\mathrm{d}\gamma}
=-\Psi(\alpha-\gamma+t)-\Psi(\beta-\gamma+t)\\&\qquad+\Psi(t-\gamma+1)+\Psi(t-k-\gamma)-\log x\end{aligned}
\]
Daher folgt für $\gamma=-k$
\SourceLeaf{225}
\[
\begin{aligned}[c]
[88]\quad &F(\alpha,\beta,\alpha+\beta+1+k,1-x)\\
&=\frac{\Pi(\alpha+\beta+k)\Pi k}{\Pi(\alpha+k)\Pi(\beta+k)}X\\
&\quad+\frac{\Pi(\alpha+\beta+k)\Pi k}{\Pi(\alpha-1)\Pi(\beta-1)\Pi(\alpha+k)\Pi(\beta+k)(-1)(-2)\ldotp\ldotp\ldotp\ldotp(-k)}\\&\qquad\Sigma\{(\log x\\
&\qquad+\Psi(\alpha+t+k)+\Psi(\beta+t+k)-\Psi(t+k+1)-\Psi t)
\\&\qquad\frac{\Pi(\alpha+t+k)\Pi(\beta+t+k)}{\Pi(t+k+1)\Pi t}x^{1+t+k}\}\\
&=\frac{\Pi(\alpha+\beta+k)\Pi k}{\Pi(\alpha+k)\Pi(\beta+k)}X
\pm\frac{\Pi(\alpha+\beta+k)x^{1+k}}{\Pi(\alpha-1)\Pi(\beta-1)\Pi(k+1)}Y
\end{aligned}
\]
Dabei gilt
\[
\begin{aligned}[c]
&Y=\{\log x+\Psi(\alpha+k)+\Psi(\beta+k)-\Psi(k+1)-\Psi(0)\}\\&\qquad F(\alpha+k+1,\beta+k+1,k+2,x)\\
&+A\frac{\alpha+k+1.\beta+k+1}{1. k+2}x\\
&+(A+B)\frac{\alpha+k+1.\alpha+k+2.\beta+k+1.\beta+k+2}{1.2. k+2. k+3}xx\\
&+(A+B+C)\\&\qquad\frac{\alpha+k+1.\alpha+k+2.\alpha+k+3.\beta+k+1.\beta+k+2.\beta+k+3}{1.2.3. k+2. k+3. k+4}x^3\\
&+\text{usw.}
\end{aligned}
\]
und
\[
\begin{aligned}[c]
A&=\frac{1}{\alpha+k+1}+\frac{1}{\beta+k+1}-\frac{1}{k+2}-1\\
B&=\frac{1}{\alpha+k+2}+\frac{1}{\beta+k+2}-\frac{1}{k+3}-\frac12\\
C&=\frac{1}{\alpha+k+3}+\frac{1}{\beta+k+3}-\frac{1}{k+4}-\frac13,\\
&\text{usw.}
\end{aligned}
\]
und das obere oder untere Zeichen ist zu nehmen, je nachdem, ob $k$ eine gerade oder eine ungerade Zahl ist.

\SixthArticle{46}
Auf diese Weise wird also $F(\alpha,\beta,\alpha+\beta+1-\gamma,1-x)$ umgeformt, wenn $\gamma$ den Wert $0$ hat oder eine negative ganze Zahl ist. Den Fall $\gamma=+1$ können wir genau ebenso behandeln; oder wir setzen kürzer in den vorangegangenen Überlegungen $k=-1$. Dann verschwindet $X$ vollständig, und wir erhalten:
\SourceLeaf{226}
\[
\begin{aligned}[c]
[89]\quad &F(\alpha,\beta,\alpha+\beta,1-x)\\
&=-\frac{\Pi(\alpha+\beta-1)}{\Pi(\alpha-1)\Pi(\beta-1)}
\\&\qquad\{\log x+\Psi(\alpha-1)+\Psi(\beta-1)-2\Psi(0)\}F(\alpha,\beta,1,x)\\
&\quad-\frac{\Pi(\alpha+\beta-1)}{\Pi(\alpha-1)\Pi(\beta-1)}
\{A\frac{\alpha.\beta}{1}x\\
&\quad+(A+B)\frac{\alpha.\alpha+1.\beta.\beta+1}{1.2.1.2}xx\\
&\quad+(A+B+C)\frac{\alpha.\alpha+1.\alpha+2.\beta.\beta+1.\beta+2}{1.2.3.1.2.3}x^3\\
&\quad+\text{usw.}\}
\end{aligned}
\]
Dabei gilt
\[
\begin{aligned}&A=\frac1\alpha+\frac1\beta-2,\\&
B=\frac1{\alpha+1}+\frac1{\beta+1}-\frac22.\\&
C=\frac1{\alpha+2}+\frac1{\beta+2}-\frac23,\\&\text{usw.}\end{aligned}
\]
So erhalten wir zum Beispiel für $\alpha=\frac12$, $\beta=\frac12$ (vergleiche Formeln 52, 71)
\[
\begin{aligned}[c]
[90]\quad &F(\tfrac12,\tfrac12,1,1-x)\\
&=-\frac1\pi\log\tfrac1{16}x. F(\tfrac12,\tfrac12,1,x)\\
&\quad-\frac1\pi\{2.\frac{1.1}{2.2}x
\\&\qquad+(2+\tfrac13)\frac{1.1.3.3}{2.2.4.4}xx
\\&\qquad+(2+\tfrac13+\tfrac2{15})\frac{1.1.3.3.5.5}{2.2.4.4.6.6}x^3\\
&\qquad+(2+\tfrac4{3.4}+\tfrac4{5.6}+\tfrac4{7.8})
\frac{1.1.3.3.5.5.7.7}{2.2.4.4.6.6.8.8}x^4+\text{usw.}\}\\
&=-\frac1\pi\{\log\tfrac1{16}xF(\tfrac12,\tfrac12,1,x)
\\&\qquad+\tfrac12x+\tfrac{21}{64}xx+\tfrac{185}{768}x^3+\tfrac{18655}{98304}x^4\\
&\qquad+\tfrac{102501}{655360}x^5+\tfrac{1394239}{10485760}x^6+\text{usw.}\}
\end{aligned}
\]
Den dritten Fall schließlich, in dem $\gamma$ eine positive ganze Zahl größer als eins ist, braucht man nicht gesondert zu behandeln, denn es gilt
\[
\begin{aligned}&F(\alpha,\beta,\alpha+\beta+1-\gamma,1-x)\\&=x^{\gamma-1}F(\alpha+1-\gamma,\beta+1-\gamma,\alpha+\beta+1-\gamma,1-x)\end{aligned}
\]
und die Umformung der Reihe $F(\alpha+1-\gamma,\beta+1-\gamma,\alpha+\beta+1-\gamma,1-x)$ für $\gamma>1$ führt von selbst auf den ersten Fall zurück.

\SixthArticle{47}
Wir gehen zu anderen Umformungen über, unter denen die Substitution $x=\frac{y}{y-1}$ die erste Stelle einnehmen soll. Daraus folgt $\mathrm{d}x=-\frac{\mathrm{d}y}{(y-1)^2}$ und somit
\[
\frac{\mathrm{d}P}{\mathrm{d}x}=-\frac{\mathrm{d}P}{\mathrm{d}y}(1-y)^2,
\quad\text{nochmaliges Differenzieren ergibt}
\]
\[
\mathrm{d}\frac{\mathrm{d}P}{\mathrm{d}x}=-(1-y)^2\mathrm{d}\frac{\mathrm{d}P}{\mathrm{d}y}
+2(1-y)\mathrm{d}P,\quad\text{somit}
\]
\[
\frac{\mathrm{dd}P}{\mathrm{d}x^2}=+(1-y)^4\frac{\mathrm{dd}P}{\mathrm{d}y^2}
-2(1-y)^3\frac{\mathrm{d}P}{\mathrm{d}y}
\]
\SourceLeaf{227}
Setzt man diese Werte ein, so geht Gleichung 80 über in
\[
\begin{aligned}&0=\alpha\beta P+(1-y)(\gamma+(\alpha+\beta-1-\gamma)y)\frac{\mathrm{d}P}{\mathrm{d}y}\\&+(1-y)(y-yy)\frac{\mathrm{dd}P}{\mathrm{d}y^2}\end{aligned}
\]
Um aber eine Gleichung von ähnlicher Form wie 80 zu erhalten, setzen wir $P=(1-y)^\mu P'$. Daraus folgt
\[
\begin{aligned}[c]
\frac{\mathrm{d}P}{\mathrm{d}y}
&=-\mu(1-y)^{\mu-1}P'+(1-y)^\mu\frac{\mathrm{d}P'}{\mathrm{d}y}\\
\frac{\mathrm{dd}P}{\mathrm{d}y^2}
&=(\mu\mu-\mu)(1-y)^{\mu-2}P'-2\mu(1-y)^{\mu-1}\frac{\mathrm{d}P'}{\mathrm{d}y}
\\&\qquad+(1-y)^\mu\frac{\mathrm{dd}P'}{\mathrm{d}y^2}
\end{aligned}
\]
Nach dem Einsetzen dieser Ausdrücke und der Division durch $(1-y)^\mu$ ergibt sich
\[
\begin{aligned}[c]
&0=P'\{\alpha\beta-\mu(\gamma+(\alpha+\beta-1-\gamma)y+y(\mu\mu-\mu)\}\\
&+\frac{\mathrm{d}P'}{\mathrm{d}y}\{\gamma+(\alpha+\beta-1-\gamma)y-2\mu y\}(1-y)\\
&+\frac{\mathrm{dd}P'}{\mathrm{d}y^2}\{y-yy\}(1-y)
\end{aligned}
\]
Bestimmen wir $\mu$ so, dass der Faktor von $P'$ durch $1-y$ teilbar wird. Dies erreicht man, indem man entweder $\mu=\alpha$ oder $\mu=\beta$ setzt. Der erste Wert verwandelt die vorangehende Gleichung in
\[
\begin{aligned}&0=\alpha(\beta-\gamma)P'+(\gamma-(\gamma+\alpha+1-\beta)y)\frac{\mathrm{d}P'}{\mathrm{d}y}\\&+(y-yy)\frac{\mathrm{dd}P'}{\mathrm{d}y^2}\end{aligned}
\]
oder
\[
\begin{aligned}&0=\alpha(\gamma-\beta)P'-(\gamma-(\gamma-\beta+\alpha+1)y)\frac{\mathrm{d}P'}{\mathrm{d}y}\\&-(y-yy)\frac{\mathrm{dd}P'}{\mathrm{d}y^2}\end{aligned}
\]
die so zu erfüllen ist, dass für $y=0$ die Gleichungen $P'=1$ und $\frac{\mathrm{d}P'}{\mathrm{d}y}=\frac{\alpha(\gamma-\beta)}{\gamma}$ gelten. Daraus folgt aber $P'=F(\alpha,\gamma-\beta,\gamma,y)$, und somit gilt
\[
\begin{aligned}&[91]\quad F(\alpha,\beta,\gamma,x)=(1-y)^\alpha F(\alpha,\gamma-\beta,\gamma,y)\\&=(1-x)^{-\alpha}F(\alpha,\gamma-\beta,\gamma,-\tfrac{x}{1-x})\end{aligned}
\]
Hätten wir für $\mu$ den anderen Wert $\beta$ genommen, so hätte sich auf genau dieselbe Weise ergeben:
\[
[92]\quad F(\alpha,\beta,\gamma,x)=(1-x)^{-\beta}F(\beta,\gamma-\alpha,\gamma,-\tfrac{x}{1-x})
\]
Diese Formel folgt auch von selbst aus der vorangehenden, wenn man lediglich die Elemente $\alpha$, $\beta$ vertauscht. Mithilfe der eben gefundenen Formel lassen sich die Werte unserer Reihen für negative Werte des vierten Elements stets auf die Werte ähnlicher Reihen für positive, zwischen $0$ und $1$ liegende Werte des vierten Elements zurückführen, denn es gilt
\[
F(\alpha,\beta,\gamma,-x)=(1+x)^{-\alpha}F(\alpha,\gamma-\beta,\gamma,\tfrac{x}{1+x})
\]
\SourceLeaf{228}
